% Source: https://artofproblemsolving.com/community/c3h1466161_practice_from_an_old_mock_not_made_by_me
% Stable identity: aops-topic-1466161-q23
The side of the hexagon equals the radius $r$. Two vertices with one vertex between them are $\sqrt{3}\,r$ apart, and opposite vertices are $2r$ apart (a diameter). There are $\dbinom{6}{4}=15$ ways to choose four vertices; equivalently, choose the two vertices that are left out.
\begin{itemize}
\item \emph{The two omitted vertices are adjacent} ($6$ ways). The quadrilateral consists of three consecutive sides and a diameter: perimeter $r+r+r+2r=5r$.
\item \emph{The omitted vertices have one vertex between them} ($6$ ways). The quadrilateral has two sides $r$ and two diagonals $\sqrt{3}\,r$: perimeter $(2+2\sqrt{3})r\approx5.46r$.
\item \emph{The omitted vertices are opposite} ($3$ ways). The quadrilateral is a rectangle with sides $r$ and $\sqrt{3}\,r$: perimeter $(2+2\sqrt{3})r$.
\end{itemize}

\begin{asy}
usepackage("amsmath");
size(9cm);
defaultpen(fontsize(10pt));
pair V(int k){return dir(60*k);}
void hexpic(pair o,int[] keep,string cap,pen p){
  draw(shift(o)*unitcircle,gray);
  for(int k=0;k<6;++k){draw(shift(o)*(V(k)--V(k+1)),gray+dashed); dot(o+V(k),gray);}
  path q=o+V(keep[0]);
  for(int i=1;i<4;++i) q=q--(o+V(keep[i]));
  draw(q--cycle,p+linewidth(0.9));
  for(int i=0;i<4;++i) dot(o+V(keep[i]));
  label(cap,o+(0,-1.35));
}
hexpic((0,0),new int[]{0,1,2,3},"$r+r+r+2r=5r$",blue);
hexpic((2.9,0),new int[]{0,1,2,4},"$2r+2\sqrt{3}\,r$",red);
hexpic((5.8,0),new int[]{0,1,3,4},"$2r+2\sqrt{3}\,r$",heavygreen);
label("$r$",0.5*(V(0)+V(1)),dir(210));
label("$2r$",0.5*(V(3)+V(0)),S);
label("\mbox{}",(min(currentpicture,user=true)+max(currentpicture,user=true))/2,black); // keeps the text after the figure black when LaTeX typesets the labels inline
\end{asy}

The perimeter is at most $5r$ exactly in the first case, so the probability is $\dfrac{6}{15}=\boxed{\dfrac{2}{5}}$.

(No choice gives a perimeter \emph{strictly} less than $5r$: the smallest possible perimeter is exactly $5r$.)
